\par\Needspace{6\baselineskip}\medskip
\begin{center}\hypertarget{native-197-article147}{}\pdfbookmark[2]{147.}{article-147}147.\end{center}
Proprie quidem etiam hic, sicuti in decem methodis supra traditis, arbitrarium
esset, quosnam valores novos pro $x$ et $y$ in hypothesi secunda supponamus,
si modo conditionibus generalibus supra explicatis non adversentur:
attamen quum manifesto pro lucro magno habendum sit, si statim a valoribus
magis exactis proficisci liceat, in methodo hacce parum prudenter ageres, si
valores secundos temere quasi adoptares, quum ex ipsa rei natura facile perspiciatur,
si valores primi ipsarum $P,Q$ levibus erroribus affecti fuerint, ipsas
$P',Q'$ valores multo exactiores exhibituras esse, siquidem motus heliocentricus
fuerit modicus. Quamobrem semper ipsas $P',Q'$ pro valoribus secundis
ipsarum $P,Q$ adoptabimus, sive $\log P',\log Q'$ pro valoribus secundis ipsarum
$x,y$, si $\log P,\log Q$ primos designare suppositi sint.

Iam in hac hypothesi secunda, ubi omnes operationes praeliminares per
formulas 1—20 exhibitae invariatae retinendae sunt, calculus prorsus simili modo
repetetur. Primo scilicet determinabitur angulus $\omega$; dein $z,r',\dfrac{n'r'}{n}\mathpunct{\raisebox{.45ex}{,}}$ $\dfrac{n'r'}{n''}$
$\zeta,r,\zeta'',r'',f',f,f''$. E differentia plus minusve considerabili inter valores novos
harum quantitatum atque primos facile aestimabitur, utrum operae pretium sit,
necne, correctionem quoque temporum propter aberrationem denuo computare:
in casu posteriori temporum intervalla, adeoque etiam quantitates $\theta$ et $\theta''$,
eaedem manebunt ut ante. Denique ex $f,r',r''$; $f'',r,r'$ temporumque intervallis
eruentur $\eta,\eta''$ atque hinc valores novi ipsarum $P',Q'$, qui plerumque ab iis,
quos hypothesis prima suppeditaverat, multo minus different, quam hi ipsi a
valoribus primis ipsarum $P,Q$. Valores secundi ipsarum $X,Y$ itaque multo
\SourceLeaf{198}minores erunt, quam primi, valoresque secundi ipsarum $P',Q'$ tamquam valores
tertii ipsarum $P,Q$ adoptabuntur, et cum his calculus denuo repetetur.
Hoc igitur modo sicuti ex hypothesi secunda numeri exactiores resultaverant,
quam ex prima, ita e tertia iterum exactiores resultabunt, quam e secunda,
possentque valores tertii ipsarum $P',Q'$ tamquam quarti ipsarum $P,Q$ adoptari,
atque sic calculos toties repeti, usque dum ad hypothesin perveniatur, in qua
$X$ et $Y$ pro evanescentibus habere liceret: sed quoties hypothesis tertia nondum
sufficiens videatur, valores ipsarum $P,Q$ in hypothesi quarta adoptandos
secundum methodum in artt. 120, 121 explicatam e tribus primis deducere
praestabit, quo pacto approximatio celerior obtinebitur, raroque opus erit, ad
hypothesin quintam progredi.
